\section{Defects and Accidents}

\begin{quotex}
Among those who receive the same teaching, each one understands and assimilates it more or less completely and profoundly according to the extent of his own intellectual possibilities; and that is why selection comes about quite naturally, without which there could be no genuine hierarchy.
\flright{\textsc{Rene Guenon}, \emph{Man and his Becoming}}
\end{quotex}


After the discussion of heredity in \emph{The Great Triad}, Rene Guenon refers the reader to the chapter, \emph{Qualifications for Initiation} in \emph{Perspectives on Initiation}. This has broader interest than the narrow question of initiation, since it exposes the topic of human inequalities. After all, if the Self is somehow complicit in choosing the conditions and circumstances of birth, then why would not every Self choose to be Healthy, Wealthy, and Wise?

First of all, the following fundamental distinction must be understood.

\begin{itemize}


\item \textbf{Person}: The Person or Self (le Soi) is the transcendent and permanent principle of the being.

\item \textbf{Individual}: An individual is a transient and contingent modification of the Being. In this discussion, we mean, by ``individual'', the human being. The Ego (le moi) is the conscious element of the person.
\end{itemize}


Each individual starts from the state of manifestation in which he is situated. This state includes many determinate conditions which are of two types:

\begin{itemize}


\item Those that are generally inherent to that state. For example, the human state includes conditions which are common to all humans.

\item Those that are particular to the individual and distinguish him from other individuals.
\end{itemize}


Most contemporary thought in the West consider the individual to be the entire the human being, and the understanding of the ``Person'' is restricted to the individual. Then the individual is then divided into a body and a soul (or mind, etc.) The former is the gross aspect and the latter contains the subtle aspects; there cannot be a duality between them, since they are all modalities of the Self. As such, they interact with each other according to the Law of Correspondences. Hence, the body may influence the mind and vice versa.

\paragraph{Obstacles to Initiation}


Guenon explains the reasons why certain infirmities and bodily defects are considered obstacles to initiation. In themselves they may seem insignificant --- for example, a stutterer, or those with asymmetries in the face or limbs, cannot become a Freemason --- yet they are outward signs of certain psychic imperfections. This follows from the correspondences between the various states of the human being. The subtle states necessarily correspond to the gross, or physical, state.

Curiously, he points out that the impediments to Masonic initiation are almost identical to the impediments to Catholic Holy Orders\footnote{\url{https://www.newadvent.org/cathen/08170a.htm}}. The point is that certain defects affect the proper performance of some rites. Physical defects, as such, are not an indication of a moral failure.

There are two types of defects:

\begin{enumerate}


\item Defects that the being exhibits from birth or that develop naturally over the course of its existence as a consequence of a certain predisposition.

\item Those that are the result of an accident.
\end{enumerate}


The first type inheres in the very nature of the being, However, since nothing can happen to a being that does not correspond in some way to an essential element of his nature, then even apparent accidental infirmities cannot be considered indifferent.

\paragraph{Law of Accidents}


\begin{quotex}
``Don't you mind,'' said Puddleglum. ``There are no accidents. Our guide is Aslan.''
\flright{\textsc{C S Lewis}, \emph{The Silver Chair}}
\end{quotex}


It is certainly difficult, in contemporary worldviews, to accept that accidents are a necessary aspect of one's being. On the contrary, the commonsense view seems to be that all qualities and conditions of human birth are arbitrary and inessential. Hence, inequality in circumstances of birth is regarded as unjust. Guenon anticipates this view when he writes:

\begin{quotex}
Some may be astonished that accidental infirmities thus correspond to something in the very nature of the being affected by them, but this is after all only a direct consequence of the real relationships the being has with the environment in which he manifests itself. Everything a being undergoes, as well as all he does, constitute a modification of that being, and must necessarily correspond to one of the possibilities in his nature, so that there can be nothing that is purely accidental if this word be understood in the sense of `extrinsic', as it commonly is.
\flright{\textsc{Rene Guenon}}
\end{quotex}


\paragraph{Victim Souls}


\begin{quotex}
Now I rejoice in my sufferings for your sake, and in my flesh I complete what is lacking in Christ's afflictions for the sake of his body.
\flright{\textsc{Colossians 1:24}}
\end{quotex}


Victim souls\footnote{\url{https://catholicstand.com/victim-souls-defeat-satan/}} deliberately agree to a life of painful physical defects and suffering.

Saint Elizabeth of the Trinity\footnote{\url{https://www.carmelitedcj.org/carmel/saints-of-carmel/160-bl-elizabeth-of-the-trinity}} is an example; in her writings, she seems to have known that her role would be to intercede for the sick and suffering.


\flrightit{Posted on 2022-02-21 by Cologero}

\begin{center}* * *\end{center}

\begin{footnotesize}
\begin{sffamily}
\texttt{De Bosit on 2022-02-22 at 00:43 said:}

The Victim Souls strike me as overzealous martyrs. Too much too fast?

Unless you're born perfect you probably have to go through hell or purgatory anyway (in life or death later), and with experience maybe it can become bearable.

Then you are a demon to demons, and can execute a meek revenge.

I like this idea more than constant pain and suffering.

\hfill

\texttt{James on 2022-02-23 at 18:11 said:}

The great American mystic Flannery O'Connor describes in her short story ``A Temple of the Holy Ghost'' a young child who is told of a hermaphrodite at a local fair who would tell the paying customers who entered the tent, ``\,`God made me thisaway and if you laugh He may strike you the same way. This is the way He wanted me to be and I ain't disputing His way. I'm showing you because I got to make the best of it. I expect you to act like ladies and gentlemen. I never done it to myself nor had a thing to do with it but I'm making the best of it. I don't dispute hit. sic' ''

The hermaphrodite's words confused the young child who goes to bed imagining the hermaphrodite talking to her saying: ``\,`Raise yourself up. A temple of the Holy Ghost. You! You are God's temple, don't you know? Don't you know? God's Spirit has a dwelling in you, don't you know?' ''

The story ends with a trip to the convent where the girls witness Exposition and Benediction of the Blessed Sacrament and the young girl remembers the hermaphrodite's words just as she sees the Blessed Sacrament being raised up: ``\,`I don't dispute hit. This is the way He wanted me to be.' ''

Thus there are some things to understand:

The fact that there is life in us is signal and proof enough that regardless of the physical state of our bodies, we are Temples of the Holy Ghost and the treatment of our self in all of its layers is our sacred duty as custodians of this temple. There is no physical ``defect'' that could nullify this inherent reality.

The association the young girl makes of the Hermaphrodite's Body and the raising of the Body of Christ is a recognition that our bodily suffering and defects is a participation in the suffering and passion of Our Lord.

The mystery of why the Hermaphrodite would disavow any choice in the matter might seem like a stumbling block, but it is not necessarily. What we can understand is a reiteration that suffering and defects are not inherently to be sought after. What the Hermaphrodite's SELF chose is not the individual nastiness of defects, but rather the salvific suffering of being united with the Lord's suffering. Just as the Lord asked His Father to let the cup pass, so, too, does the Hermaphrodite understand that suffering is indeed undesirable. It is disordered. Contrary to O'Connor being read as someone who might refuse the participation of the will, I believe it is more a careful reaffirmation that suffering is an evil that will be wiped away and that the soul properly does not choose evil, but it can, however, surrender to the work of salvation.

Funny enough, at the end of the story, the fair is shut down because local Protestant preachers came around had the police close the fair. There is no solution to suffering in the Protestant mindset--either bodily or circumstantially. It is why the corpus on the crucifix is a scandal to many southerners. Not even counting that we eat His flesh and drink His blood.

As far as my meditations on our existence is concerned, I tend to think that we are here for one of two reasons. We chose this most bleak time of our history because we are either one of the thieves crucified along with the Lord, or we are those who have the moral courage to undergo intense spiritual training (c.f. the unity of the Lion and the Angel ``to dare'' in the Wheel of Fortune). In other words, I can choose right now to turn to the Lord who is suffering with me and tell Him that it is not at all unfair that I am subject to a godless society for I myself deserve this curse for my crimes and that all I ask is that He remembers me when He enters into His Kingdom. Otherwise, I am a daring soul who has decided to participate in the age of catacombic heroism. I have not yet penetrated into which one I am yet. I tend to think with all my sins I am the thief and I tend to think with all of my faith that I am the one who dares with moral courage. Perhaps everyone who is born in this time is a little of both.

\hfill

\texttt{Squirrel on 2022-03-10 at 12:40 said:}

Personal infamy such as being an actor or working in a slaughterhouse was also a bar (though almost always overlooked) to taking the Chinese standardized test, although eunuchs and the crippled were eligible since they wanted as wide a base as possible.

\end{sffamily}
\end{footnotesize}

